
\documentclass[a4paper,11pt]{article}
\usepackage[a4paper, margin=8em]{geometry}

% usa i pacchetti per la scrittura in italiano
\usepackage[french,italian]{babel}
\usepackage[T1]{fontenc}
\usepackage[utf8]{inputenc}
\frenchspacing 

% usa i pacchetti per la formattazione matematica
\usepackage{amsmath, amssymb, amsthm, amsfonts}

% usa altri pacchetti
\usepackage{gensymb}
\usepackage{hyperref}
\usepackage{standalone}

% imposta il titolo
\title{Appunti Comunicazioni Numeriche}
\author{Luca Seggiani}
\date{2026}

% imposta lo stile
% usa helvetica
\usepackage[scaled]{helvet}
% usa palatino
\usepackage{palatino}
% usa un font monospazio guardabile
\usepackage{lmodern}

\renewcommand{\rmdefault}{ppl}
\renewcommand{\sfdefault}{phv}
\renewcommand{\ttdefault}{lmtt}

% disponi il titolo
\makeatletter
\renewcommand{\maketitle} {
	\begin{center} 
		\begin{minipage}[t]{.8\textwidth}
			\textsf{\huge\bfseries \@title} 
		\end{minipage}%
		\begin{minipage}[t]{.2\textwidth}
			\raggedleft \vspace{-1.65em}
			\textsf{\small \@author} \vfill
			\textsf{\small \@date}
		\end{minipage}
		\par
	\end{center}

	\thispagestyle{empty}
	\pagestyle{fancy}
}
\makeatother

% disponi teoremi
\usepackage{tcolorbox}
\newtcolorbox[auto counter, number within=section]{theorem}[2][]{%
	colback=blue!10, 
	colframe=blue!40!black, 
	sharp corners=northwest,
	fonttitle=\sffamily\bfseries, 
	title=~\thetcbcounter: #2, 
	#1
}

% disponi definizioni
\newtcolorbox[auto counter, number within=section]{definition}[2][]{%
	colback=red!10,
	colframe=red!40!black,
	sharp corners=northwest,
	fonttitle=\sffamily\bfseries,
	title=~\thetcbcounter: #2,
	#1
}

% U.D.M
\newcommand{\amp}{\ensuremath{\mathrm{A}}}
\newcommand{\volt}{\ensuremath{\mathrm{V}}}
\newcommand{\meter}{\ensuremath{\mathrm{m}}}
\newcommand{\second}{\ensuremath{\mathrm{s}}}
\newcommand{\farad}{\ensuremath{\mathrm{F}}}
\newcommand{\henry}{\ensuremath{\mathrm{H}}}
\newcommand{\siemens}{\ensuremath{\mathrm{S}}}

% circuiti
\usepackage{circuitikz}
\usetikzlibrary{babel}

% disegni
\usepackage{pgfplots}
\pgfplotsset{width=10cm,compat=1.9}

% disponi codice
\usepackage{listings}
\usepackage[table]{xcolor}

\lstdefinestyle{codestyle}{
		backgroundcolor=\color{black!5}, 
		commentstyle=\color{codegreen},
		keywordstyle=\bfseries\color{magenta},
		numberstyle=\sffamily\tiny\color{black!60},
		stringstyle=\color{green!50!black},
		basicstyle=\ttfamily\footnotesize,
		breakatwhitespace=false,         
		breaklines=true,                 
		captionpos=b,                    
		keepspaces=true,                 
		numbers=left,                    
		numbersep=5pt,                  
		showspaces=false,                
		showstringspaces=false,
		showtabs=false,                  
		tabsize=2
}

\lstdefinestyle{shellstyle}{
		backgroundcolor=\color{black!5}, 
		basicstyle=\ttfamily\footnotesize\color{black}, 
		commentstyle=\color{black}, 
		keywordstyle=\color{black},
		numberstyle=\color{black!5},
		stringstyle=\color{black}, 
		showspaces=false,
		showstringspaces=false, 
		showtabs=false, 
		tabsize=2, 
		numbers=none, 
		breaklines=true
}

\lstdefinelanguage{javascript}{
	keywords={typeof, new, true, false, catch, function, return, null, catch, switch, var, if, in, while, do, else, case, break},
	keywordstyle=\color{blue}\bfseries,
	ndkeywords={class, export, boolean, throw, implements, import, this},
	ndkeywordstyle=\color{darkgray}\bfseries,
	identifierstyle=\color{black},
	sensitive=false,
	comment=[l]{//},
	morecomment=[s]{/*}{*/},
	commentstyle=\color{purple}\ttfamily,
	stringstyle=\color{red}\ttfamily,
	morestring=[b]',
	morestring=[b]"
}

% disponi sezioni
\usepackage{titlesec}

\titleformat{\section}
	{\sffamily\Large\bfseries} 
	{\thesection}{1em}{} 
\titleformat{\subsection}
	{\sffamily\large\bfseries}   
	{\thesubsection}{1em}{} 
\titleformat{\subsubsection}
	{\sffamily\normalsize\bfseries} 
	{\thesubsubsection}{1em}{}

% disponi alberi
\usepackage{forest}

\forestset{
	rectstyle/.style={
		for tree={rectangle,draw,font=\large\sffamily}
	},
	roundstyle/.style={
		for tree={circle,draw,font=\large}
	}
}

% disponi algoritmi
\usepackage{algorithm}
\usepackage{algorithmic}
\makeatletter
\renewcommand{\ALG@name}{Algoritmo}
\makeatother

% disponi numeri di pagina
\usepackage{fancyhdr}
\fancyhf{} 
\fancyfoot[L]{\sffamily{\thepage}}

\makeatletter
\fancyhead[L]{\raisebox{1ex}[0pt][0pt]{\sffamily{\@title \ \@date}}} 
\fancyhead[R]{\raisebox{1ex}[0pt][0pt]{\sffamily{\@author}}}
\makeatother

\begin{document}
% sezione (data)
\section{Lezione del 31-03-26}

% stili pagina
\thispagestyle{empty}
\pagestyle{fancy}

% testo
Avevamo visto nella lezione le variabili aleatorie continue, cioè quelle che assumevano un numero infinito non numerabile di valori. Quindi, avevamo visto come la definizione di \textit{densità} di probabilità veniva definita come derivata:
$$
f_X(x) = \frac{dF_X (x)}{dx}
$$

\subsubsection{Esempi di variabili aleatorie continue}
Vediamo quindi alcuni esempi di variabili aleatorie continue:
\begin{itemize}
	\item \textbf{\textsf{Variabili aleatorie uniformi}} \\
Una variabile aleatoria è detta \textit{uniforme} nell'intervallo $[a, b]$, se la sua densità di probabilità è costante in tale intervallo:
$$
f_X(x) = 
\begin{cases}
	\frac{1}{b-a}, \quad a \leq x \leq b \\
	0, \quad \text{altrove}
\end{cases}, \quad
F_X(x) = 
\begin{cases}
	0, \quad x < a \\
	\frac{x - a}{b-a}, \quad a \leq x \leq b \\
	1, \quad x > b 
\end{cases}
$$
Come abbiamo detto, $X$ è uniforme e si dice:
$$
X \in \mathrm{Uniforme}(a, b)
$$

	\item \textbf{\textsf{Variabili aleatorie esponenziali}} \\
Vediamo un altro esempio che tornerà spesso di probabilità di densità, cioè la variabile aleatoria \textit{esponenziale}. Noteremo che questa ha la proprietà di \textit{assenza di memoria}. Ha un unico parametro $\lambda$, ed è definita come:
$$
f_X(x) = \frac{1}{\lambda} e^{-\frac{x}{\lambda}} u(x)
, \quad
F_X(x) = \left( 1 - e^{-\frac{x}{\lambda}} \right) u(x)
$$
dove $u$ è il gradino di Heaviside.

La distribuzione di probabilità si calcola come sempre svolgendo l'integrale:
$$
F_X(x) =
\int_{-\infty}^x \frac{1}{\lambda} e^{-\frac{\alpha}{x}} u(\alpha) \, d\alpha =
\begin{cases}
	0, \quad x < 0 \\
	\frac{1}{\lambda} \int_0^x e^{-\frac{\alpha}{x}} \, d\alpha, \quad x \geq 0
\end{cases} 
= \left( 1 - e^{-\frac{x}{\lambda}} \right) u(x)
$$
dove si reintroduce il gradino $u$ alla fine della derivazione. Come abbiamo detto, $X$ è esponenziale e si dice:
$$
X \in \mathrm{Esponenziale}(\lambda)
$$
	
\item \textbf{\textsf{Variabili aleatorie gaussiane}} \\
Veniamo quindi al tipo più celebre di variabile aleatoria, cioè quella \textit{gaussiana} o \textit{normale}. Questa è indicata da due parametri caratteristici:
\begin{itemize}
	\item Il \textit{valor medio} $\eta$. Corrisponde al \textit{massimo} della densità di probabilità $f_X(x)$;
	\item La \textit{varianza} $\sigma^2$, cioè un indice di variazione dal valor medio della variabile aleatoria. Corrisponde al \textit{flesso} $\eta \pm \sigma$ della densità di probabilità $f_X(x)$.
\end{itemize}
Viene quindi definita come:
$$
f_X(x) = \frac{1}{\sigma \sqrt{2 \pi}} e^{-\frac{1}{2} \left( \frac{x - \eta}{\sigma} \right)^2} = \frac{1}{\sqrt{2 \pi \sigma^2}} e^{- \frac{(x - \eta)}{2\sigma^2}}
$$
Se una $X$ è normale o gaussiana scriviamo:
$$
X \in \mathcal{N}(\eta, \sigma^2)
$$

Notiamo che per una variabile aleatoria gaussiana non è possibile calcolare in forma chiusa la distribuzione di probabilità $F_X(x)$. Per questo l'unico modo che avremo per ricavarla sarà per via numerica o attraverso tabelle.

\par\smallskip

Talvolta definiamo la variabile aleatoria gaussiana (o normale) \textit{standard} come:
$$
f_z (z) = \frac{1}{\sqrt{2 \pi}} e^{-\frac{z^2}{2}}
$$
presi valor medio e varianza come:
$$
\begin{cases}
	\eta = 0 \\
	\sigma = 1
\end{cases}
$$
Chiamiamo quindi la funzione di densità $\Phi$, indicata come:
$$
\Phi(z) = F_Z(z) = \frac{1}{\sqrt{2 \pi}} \int_{-\infty}^z e^{-\frac{\alpha^2}{2}} \, d\alpha
$$
Prendiamo quindi la speculare $Q$ come la \textit{funzione degli errori}:
$$
Q = 1 - \Phi(z)
$$

\par\smallskip

Se $X \in \mathcal{N}(\eta, \sigma^2)$, la funzione di distribuzione si ricava da quella della normale standard mediante la seguente relazione:
$$
F_X(x) = \frac{1}{\sigma \sqrt{2 \pi}} \int_{-\infty}^x e^{-\frac{(\alpha - \eta)^2}{2 \sigma^2}} \, d\alpha
$$
quindi ponendo il cambio di variabile in $y$:
$$
y = \frac{\alpha - \eta}{\sigma}, \quad dy = \frac{d\alpha}{\sigma}
$$
per cui ritroviamo la stessa forma della funzione $\Phi$:
$$
F_X(x) = \frac{1}{\sqrt{2 \pi}} \int_{-\infty}^{\frac{x - \alpha}{\sigma}} e^{-\frac{y^2}{2}} \, dy = \Phi \left( \frac{x - \eta}{\sigma} \right) = 1 - Q \left( \frac{x - \eta}{\sigma} \right)
$$
Quindi si ha che:
$$
P(X > \lambda) = 1 - F_X(\lambda) = Q \left( \frac{\lambda - \eta}{\sigma} \right)
$$
cioè la probabilità che la variabile aleatoria gaussiana $X$ si trovi sopra $\lambda$ è uguale alla funzione degli errori $Q$ valutata nel punto $y$ con $\alpha = \lambda$. Ricordiamo ancora che $Q$ si ricava da tabelle o in maniera numerica.

\par\smallskip

In generale, data una variabile aleatoria gaussiana $X \in \mathcal{N}(\eta, \sigma^2)$, la probabilità associata ad un qualunque evento di interesse può essere calcolata mediante la funzione $\Phi$, o equivalentemente mediante la funzione $Q$:
$$
P(a < x \leq b) = \int_a^b f_X(\alpha) \, d\alpha = F_X(b) - F_X(a) 
= \Phi\left(\frac{b - \eta}{\sigma}\right) - \Phi\left(\frac{a - \eta}{\sigma}\right)
= Q\left(\frac{a - \eta}{\sigma}\right) - Q\left(\frac{b - \eta}{\sigma}\right)
$$
in quanto per simmetria vale che:
$$
Q(z) = 1 - Q(-z), \quad \Phi(z) = 1 - \Phi(-z)
$$

\end{itemize}

\subsubsection{Densità di probabilità di variabili aleatorie discrete}
Per le variabili aleatorie discrete avevamo definito in 8.3 una funzione \textit{massa} $p_X(x)$ di probabilità. Risulta possibile anche scrivere una funzione \textit{densità di probabilità} sfruttando la delta di Dirac (introdotta per trattare le trasformate di Fourier generalizzate):
$$
F_X(x) = \sum_i p_X (x_i) u (x - x_i) \implies f_X(x) = \sum_i p_X(x_i) \delta(x - x_i)
$$

Questo risulta immediato dalla definizione di densità di probabilità come:
$$
f_X(x) = \frac{d F_X(x) }{dx}
$$
e dal fatto che la distribuzione di probabilità per una variabile aleatoria discreta era definita in funzione del gradino $u$:
$$
F_X(x) = P(\{ X \leq x \}) = \sum_i p_X (x_i) u (x - x_i)
$$
Quindi, noto il fatto che la derivata del gradino è la delta di Dirac:
$$
\frac{d}{dx} u(x) = \delta(x)
$$
il risultato ottenuto è immediato.

Questo ha un risultato anche per le probabilità di intervalli:
$$
P(a < X \leq b) = \int_a^b f_X(x) \, dx = \sum_i p_X(x_i) \int_a^b \delta(x - x_i) \, dx
$$
Se $a$ e $b$ non coincidono con nessuno degli $x_i$, si ottiene semplicemente la somma delle probabilità associate ai valori assunti dalla variabile aleatoria all'interno dell'intervallo $a$, $b$. Si deve fare attenzione, ricordiamo, al caso in cui uno degli estremi (facciamo $b$) coincide con un valore della variabile aleatoria. In tal caso i due integrali:
$$
\int_a^{b^+} f_X(x) \, dx \neq \int_a^{b^-} f_X(x) \, dx
$$
Nel primo caso prendiamo la densità di probabilità in $b$, nel secondo no. 

\subsection{Funzioni di distribuzione di variabili aleatorie miste}
Diamo la definizione di variabile aleatoria \textbf{mista}:
\begin{definition}{Variabile aleatoria mista}
	Una variabile aleatoria si dice mista se la relativa funzione di distribuzione $F_X(x)$ è discontinua, ma non costante a tratti.
\end{definition}

In tal caso, in un intervallo continuo, può assumere valori discreti con probabilita non nulla. In particolare, un salto in $x_i$ rappresenta la probabilità dell'evento $X = x_i$. 

\subsection{Trasformazioni di variabile aleatoria}
Definiamo il concetto di \textbf{trasformazione} di una variabile aleatoria:
\begin{definition}{Trasformazione di variabile aleatoria}
Una trasformazione di variabile aleatoria è una funzione $g$ che porta una variabile aleatoria $X$ in una seconda variabile aleatoria $Y$:
$$
Y = g(X)
$$
\end{definition}

In questo una trasformazione di variabile aleatoria è una funzione reale di variabile reale, non necessariamente iniettiva. La variabile reale nel dominio comprende tutti possibili valori di $X$, e quindi il codominio corisponde a tutti i possibili valori di $Y$ dove quegli $X$ mappano. 

\subsubsection{Metodo della funzione distribuzione}

Per calcolare la probabilità (e quindi la distribuzione di probabilità) di una variabile aleatoria trasformata dobbiamo riportare l'intervallo di probabilità in $Y$ ad un intervallo $I(y)$ nella $X$ originale.
$$
F_Y(y) = P(Y \leq y) = P(X \in I(y)), \quad \text{dove} \ I(y) = \{ x: g(x) \leq y \}
$$
Da questo possiamo calcolare la distribuzione di $Y$ a partire dalla densità di $X$ come:
$$
F_Y(y) = \int_{I(y)} f_X(x) \, dx
$$

Una volta ricavata la funzione di distribuzione è possibile ricavare la densità di probabilità in $Y$, come sempre, derivando:
$$
f_Y(y) = \frac{d F_Y(y)}{dy}
$$

Questo metodo, chiamato \textit{metodo della funzione distribuzione}, è generale, cioè si applica a variabili aleatorie continue, discrete e miste. Esistono però alcuni casi dove è possibile ricavare la densità di probabilità della variabile aleatoria trasformata in maniera più semplice.

\par\smallskip

Se $Y$ è una variabile aleatoria discreta che assume i valori $y_1, y_2, ...$ finiti o infiniti numerabili, è spesso conveniente determinare direttamente la massa di probabilità $P(Y = y_k)$. Se anche la $X$ è discreta, l'evento $\{g(X) = y_k\}$ non sarà altro che l'unione di tutti gli eventi $\{X = x_i\}$ tali che $g(x_i) = y_k$.

\end{document}
